\subsection{A pointed integral-center criterion}
\begin{proposition}
Let $\mathcal C$ be pseudo-unitary. Suppose its simple dimensions are
$a_i+b_i\sqrt q$, with $b_i\geq0$, and $b_i=0$ precisely for invertible
simples. If the tensor unit is the only integral-dimensional simple
constituent of $I(\mathbf1)$, then the integral-dimensional subcategory
of $Z(\mathcal C)$ is pointed, and the forgetful functor is injective
on its invertible objects.
\end{proposition}
\begin{proof}
An integral-dimensional central simple restricts to the pointed source
subcategory, since the irrational coefficients of source dimensions are
nonnegative. By adjunction, such a simple can contain the source unit
after forgetting only if it is itself the unit. For an integral central
simple $Z$, every constituent of $ZZ^*$ is integral: its restriction
belongs to the pointed source subcategory. Hence
\[
 1=[\mathbf1:ZZ^*]=[\mathbf1:F(ZZ^*)]
   =\sum_{g\in G(\mathcal C)}[g:F(Z)]^2.
\]
Thus $F(Z)$ is a single invertible and $Z$ is invertible. Two such
invertibles with the same restriction have quotient restricting to the
unit, so their quotient is the unit.
\end{proof}
For commutative source rings, the hypothesis on $I(\mathbf1)$ is checked
on the complete distinguished dimension list $D/f_E$. The criterion
removes every integral high restriction vector other than a single
source invertible. Its exact quadratic hypotheses and every deleted
column are recorded for $R8\_3\_1768$ and $R8\_3\_1795$.
